\section{The non-normal fibre: exact claim under investigation}

\paragraph{Status convention.}
This notebook does not assert that the CDP theorem is false. It studies one
specific inference in the printed proof and the conditional application of a
conductor-supported differential to the candidate fibre $W$. A defect in the
printed inference, an alternative proof of the published theorem, and failure
of the candidate construction are logically distinct possibilities.

\subsection{The local conductor proposition}

\begin{proposition}[conductor-supported ambient differential]
\label[proposition]{prop:conductor-ambient}
\Verified\ Let $g:\mathcal X\to\Delta$ be proper, with $\mathcal X$ a
smooth complex threefold, and let $S=g^{-1}(0)$ be a reduced
normal-crossings surface. Write $\eta:\widetilde S\to S$ for the
normalization, $D\subset S$ for the reduced singular locus, and
$\widetilde D=\eta^{-1}(D)$. Let $A$ be a line bundle on $S$ such that
$\eta^*A\simeq\mathcal O_{\widetilde S}$, and choose a nowhere-vanishing
section $e$ of $\eta^*A$. Then there is a nonzero section
\[
 s\in H^0\!\left(S,\Omega^1_{\mathcal X}|_S\otimes A\right)
\]
whose image in $\Omega^1_S\otimes A$ is supported on $D$, hence is
torsion, and whose image in
\[
 \widetilde\Omega^1_S\otimes A,
 \qquad
 \widetilde\Omega^1_S:=\Omega^1_S/\mathcal T(\Omega^1_S),
\]
is zero. If $S$ is connected and $A\not\simeq\mathcal O_S$, then
$H^0(S,N^*_{S/\mathcal X}\otimes A)=H^0(S,A)=0$, so $s$ does not come
from the conormal term.
\end{proposition}

\begin{proof}
Locally $g=ux$, $uxy$, or $uxyz$, where $u$ is a unit. The conormal map is
generated by $dg|_S$. At a double point,
\[
 dg|_S=u(y\,dx+x\,dy),
\]
and at a triple point it has coefficients in $(xy,xz,yz)$. Thus $dg|_S$
is nonzero away from $D$ but restricts to zero on $D$; equivalently, it
lies in the conductor ideal times $\Omega^1_{\mathcal X}|_S$.

The normalization--conductor square gives
\[
0\longrightarrow\mathcal O_S\longrightarrow
 \eta_*\mathcal O_{\widetilde S}\oplus\mathcal O_D
 \longrightarrow\eta_*\mathcal O_{\widetilde D}\longrightarrow0.
\]
After tensoring with the locally free sheaf
$E=\Omega^1_{\mathcal X}|_S\otimes A$, the compatible pair
\[
 \bigl(\eta^*(dg|_S)\otimes e,\,0\bigr)
\]
descends uniquely to $s$. It is nonzero off $D$. Its intrinsic differential
vanishes off $D$, so its image in $\Omega^1_S\otimes A$ is
conductor-supported torsion. For nontrivial $A$, a global section of $A$
pulls back to constants on the compact connected components of
$\widetilde S$. Compatibility across the conductor and connectedness of
$S$ force either all constants to vanish or all to be nonzero; the latter
would trivialize $A$. Hence $H^0(S,A)=0$.
\end{proof}

Consequently the implication
\[
\left[
 H^0(S,N^*_{S/\mathcal X}\otimes A)=0,\quad
 H^0(S,\widetilde\Omega^1_S\otimes A)=0
\right]
\Longrightarrow
H^0(S,\Omega^1_{\mathcal X}|_S\otimes A)=0
\]
is not a formal consequence of the conormal sequence for a non-normal
normal-crossings surface. The official CDP article prints precisely this
reduction on p.~692 \cite[p.~692]{CDP2020}. This observation concerns that inference, not by itself
the truth value of the theorem.

\begin{lemma}[pullback triviality on a $\mathrm{dP}_6$ normalization]
\label[lemma]{lem:dp6-pullback}
\Verified\ Let $i:W\hookrightarrow X$ be a closed complex subspace with
$H^2(X;\mathbb Z)=0$, and suppose the normalization
$\eta:\widetilde W\to W$ satisfies
$\widetilde W\simeq\mathrm{dP}_6$.  Then, for every
$B\in\operatorname{Pic}(X)$,
\[
 \eta^*i^*B\simeq\mathcal O_{\widetilde W}.
\]
\end{lemma}

\begin{proof}
The group containing $c_1(B)$ is zero, so
$c_1(\eta^*i^*B)=0$.  Since
$H^1(\mathrm{dP}_6,\mathcal O)=0$, the exponential sequence makes
\[
 c_1:\operatorname{Pic}(\mathrm{dP}_6)
 \longrightarrow H^2(\mathrm{dP}_6;\mathbb Z)
\]
injective.  Hence the pullback line bundle is trivial.
\end{proof}

\subsection{Candidate geometry and the topology-free
\texorpdfstring{$\operatorname{Pic}^0$}{Pic0} route}

\begin{proposition}[analytic geometry of the reconstructed cusp fibre]
\label[proposition]{prop:candidate-geometry}
\Verified\ The reconstructed source produces a connected compact smooth
complex threefold $X$, a proper surjective holomorphic map
$f:X\to B\simeq\mathbb P^1$, and a full scheme-theoretic fibre
$W=f^{-1}(p_0)$ with the following properties:
\begin{enumerate}[label=(\roman*)]
\item $W$ is a connected reduced normal-crossings Cartier surface, with
local equations $z_0$, $z_0z_1$, and $z_0z_1z_2$ at its one-, two-, and
three-branch points;
\item its normalization is $\eta:\mathrm{dP}_6\to W$;
\item $W$ has an open neighbourhood $N_0\subset X$ on which
$\omega_X$ is trivial.
\end{enumerate}
\end{proposition}

\begin{proof}
The period construction in Theorem~3.4 of the frozen transcription supplies
the admissible family and the cusp-holomorphic matrix $C(t)$.  In source
pp.~26--28 (lines 1331--1446), Theorem~4.5 proves that the toric action is
free and properly discontinuous, that its quotient $N_0$ is a connected
Hausdorff smooth threefold, that $f_0:N_0\to\Delta$ is proper, and that its
full fibre is reduced with exactly the three displayed local equations.  The
same proof constructs an invariant nowhere-vanishing holomorphic three-form
on $N_0$.

Proposition~4.6, source p.~29 (lines 1453--1465), proves that the finite map
$\mathrm{dP}_6\to W$ is surjective and generically injective and hence is
the normalization.  Proposition~4.7(i)--(iii) and Theorem~5.4 give the
holomorphic overlap maps.  Finally Theorem~6.2, source pp.~36--37
(lines 1811--1834), checks the quotient atlas, Hausdorffness, properness,
compactness, and connectedness, and identifies
$f^{-1}(D_0)=N_0$ and $f^{-1}(p_0)=W$ as a reduced divisor.  Because the
embedding $N_0\hookrightarrow X$ is open and holomorphic, the constructed
trivialization of $K_{N_0}$ is a trivialization of $\omega_X|_{N_0}$.

These source gates were rechecked independently in the present workbench;
no topology, Mayer--Vietoris, or Leray assertion is used here.
\end{proof}

\begin{proposition}[$\operatorname{Pic}^0$ pullback without global topology]
\label[proposition]{prop:pic0-pullback}
\Verified\ Let $X$ be a compact complex manifold, let
$i:W\hookrightarrow X$ have normalization
$\eta:\widetilde W\to W$ with $\widetilde W\simeq\mathrm{dP}_6$, and
suppose $\omega_X$ is trivial on a neighbourhood of $W$.  For every
$L\in\operatorname{Pic}^0(X)$, put
\[
 A_L=(L^\vee\otimes\omega_X)|_W.
\]
Then
\[
 \eta^*A_L\simeq\mathcal O_{\widetilde W}.
\]
\end{proposition}

\begin{proof}
Every $L\in\operatorname{Pic}^0(X)$ has $c_1(L)=0$.  Therefore
$c_1(\eta^*i^*L)=0$.  Since
$H^1(\mathrm{dP}_6,\mathcal O)=0$, the exponential sequence makes the
Chern-class map on $\operatorname{Pic}(\mathrm{dP}_6)$ injective, so
$\eta^*i^*L$ is trivial.  The neighbourhood trivialization of $\omega_X$
also makes $\eta^*i^*\omega_X$ trivial, proving the formula for $A_L$.
\end{proof}

For the stronger all-line-bundle assertion, define
\[
 r:=\eta^*i^*:\operatorname{Pic}(X)\longrightarrow
 \operatorname{Pic}(\mathrm{dP}_6).
\]
The preceding proof gives $r|_{\operatorname{Pic}^0(X)}=0$ and the local
canonical trivialization gives $r(\omega_X)=0$.  Hence, for arbitrary $L$,
\[
 \eta^*A_L\simeq r(L)^\vee.
\]
Thus the present conductor method extends from $\operatorname{Pic}^0(X)$
to every $L\in\operatorname{Pic}(X)$ exactly when $r=0$.  The still
conditional equality $H^2(X;\mathbb Z)=0$ is one sufficient route to $r=0$
by Lemma~\ref{lem:dp6-pullback}, but it is no longer needed for the
$\operatorname{Pic}^0$ statement relevant to the generic twists in CDP20.

\begin{proposition}[candidate conductor application on $\operatorname{Pic}^0$]
\label[proposition]{prop:candidate-conductor}
\Verified\ For every $L\in\operatorname{Pic}^0(X)$, putting
\[
 \mathcal F_L=T_X\otimes L,
 \qquad
 A_L=(L^\vee\otimes\omega_X)|_W,
\]
one has
\[
 H^0(W,\Omega_X^1|_W\otimes A_L)\ne0.
\]
\end{proposition}

\begin{proof}
Proposition~\ref{prop:candidate-geometry} supplies the proper smooth family,
the full reduced normal-crossings fibre and its $\mathrm{dP}_6$
normalization.  Proposition~\ref{prop:pic0-pullback} gives
$\eta^*A_L\simeq\mathcal O_{\mathrm{dP}_6}$.  The verified conductor
proposition then supplies the nonzero section.
\end{proof}

\begin{corollary}[conditional extension to all line bundles]
\label[corollary]{cor:candidate-conductor-all}
\Conditional\ If $r=0$, the conclusion of
Proposition~\ref{prop:candidate-conductor} holds for every
$L\in\operatorname{Pic}(X)$.
\end{corollary}

\subsection{The exact duality implication required}

\begin{proposition}[absolute duality on a compact Cartier fibre]
\label[proposition]{prop:cartier-fibre-duality}
\Verified\ Let $f:X\to B$ be a holomorphic map from a smooth complex
threefold to a smooth complex curve.  Suppose that
$i:W=f^{-1}(b)\hookrightarrow X$ is compact, is the full scheme-theoretic
fibre, and is an effective Cartier divisor.  For every line bundle $L$ on
$X$, put
\[
 \mathcal F_L=T_X\otimes L,
 \qquad
 A_L=(L^\vee\otimes\omega_X)|_W.
\]
Then $W$ is a pure two-dimensional Gorenstein complex space and there is a
functorial perfect pairing
\[
 H^2(W,i^*\mathcal F_L)\times
 H^0(W,\Omega_X^1|_W\otimes A_L)\longrightarrow\mathbb C.
\]
Equivalently,
\[
 H^2(W,i^*\mathcal F_L)^\vee
 \simeq H^0(W,\Omega_X^1|_W\otimes A_L).
\]
No normality hypothesis on $W$ is required.
\end{proposition}

\begin{proof}
Because $W$ is an effective Cartier divisor in the smooth threefold $X$,
each local ring of $W$ is a two-dimensional hypersurface quotient of a
regular local ring.  Thus $W$ is Cohen--Macaulay and Gorenstein.  If
$K_W^\bullet$ denotes the analytic dualizing complex, the Gorenstein
criterion and complete-intersection computation of Ramis--Ruget give
\[
 K_W^\bullet\simeq\omega_W[2]
\]
for an invertible dualizing sheaf $\omega_W$
\cite[Proposition~2 and the complete-intersection calculation,
pp.~90--91]{RamisRuget1970}.

Cartier adjunction follows by applying
$\mathcal Hom_X(-,\omega_X)$ to
\[
0\longrightarrow\mathcal O_X(-W)\longrightarrow\mathcal O_X
 \longrightarrow i_*\mathcal O_W\longrightarrow0:
\]
\[
 \omega_W=\mathcal Ext_X^1(i_*\mathcal O_W,\omega_X)
 \simeq(\omega_X\otimes\mathcal O_X(W))|_W.
\]
Since $W=f^*(b)$ as a scheme-theoretic Cartier divisor,
$\mathcal O_X(W)\simeq f^*\mathcal O_B(b)$.  Its restriction to $W$ is the
pullback of the one-dimensional vector space $\mathcal O_B(b)|_b$, so a
choice of nonzero vector gives
\[
 \mathcal O_X(W)|_W\simeq\mathcal O_W,
 \qquad \omega_W\simeq\omega_X|_W.
\]
The choice changes the final pairing by a nonzero scalar only.

Ramis--Ruget duality pairs, for coherent $\mathcal G$,
\[
 H_c^p(W,\mathcal G)
 \quad\hbox{with}\quad
 \operatorname{Ext}^{-p}_W(\mathcal G,K_W^\bullet)
\]
\cite[Theorem~2, pp.~82--83]{RamisRuget1970}.  Compactness gives
$H_c^p(W,\mathcal G)=H^p(W,\mathcal G)$.  With
$p=2$ and $\mathcal G=i^*\mathcal F_L$, the shift above gives
\[
 \operatorname{Ext}^{-2}_W(\mathcal G,K_W^\bullet)
 \simeq\operatorname{Hom}_W(\mathcal G,\omega_W).
\]
The coherent cohomology groups are finite-dimensional by Grauert finiteness
\cite{Grauert1960}; in particular their natural topologies are separated,
so the strong duality of Ramis--Ruget is an ordinary algebraic perfect
duality.  Since $\mathcal G$ is locally free,
\[
\begin{aligned}
 \operatorname{Hom}_W(\mathcal G,\omega_W)
 &=H^0(W,\mathcal G^\vee\otimes\omega_W)\\
 &\simeq
 H^0(W,\Omega_X^1|_W\otimes L^\vee|_W\otimes\omega_X|_W)\\
 &=H^0(W,\Omega_X^1|_W\otimes A_L).
\end{aligned}
\]
The pairing is cup product followed by evaluation and the trace
$H^2(W,\omega_W)\to\mathbb C$.  Precomposition on the Hom term shows the
required contravariance in $\mathcal G$.
\end{proof}

\begin{corollary}[the duality bridge used in Section~10]
\label[corollary]{cor:cartier-fibre-h2}
Under the hypotheses of the proposition,
\[
 H^0(W,\Omega_X^1|_W\otimes A_L)\ne0
 \quad\Longrightarrow\quad
 H^2(W,i^*\mathcal F_L)\ne0.
\]
\end{corollary}

\begin{corollary}[candidate fibrewise nonvanishing]
\label[corollary]{cor:candidate-fibrewise-h2}
\Verified\ For every $L\in\operatorname{Pic}^0(X)$,
\[
 H^2(W,i^*\mathcal F_L)\ne0.
\]
The duality isomorphism itself holds for every $L\in\operatorname{Pic}(X)$.
If $r=0$, the displayed nonvanishing also holds for every such $L$.
\end{corollary}

\begin{proof}
Proposition~\ref{prop:candidate-geometry} verifies that $W$ is the compact
full Cartier fibre in the smooth threefold, so the abstract duality theorem
applies.  Proposition~\ref{prop:candidate-conductor} supplies its nonzero
dual section for $L\in\operatorname{Pic}^0(X)$.  The final assertion uses
Corollary~\ref{cor:candidate-conductor-all}.
\end{proof}

\subsection{The exact top-degree base-change implication required}

\begin{proposition}[top-degree analytic base change]
\label[proposition]{prop:top-degree-base-change}
\Verified\ Let $f:X\to B$ be a proper morphism of complex spaces,
$\mathcal F$ a coherent sheaf flat over $B$ near $c\in B$, and
$X_c$ the analytic, hence full scheme-theoretic, fibre.  If
$H^3(X_c,\mathcal F|_{X_c})=0$, then the canonical map
\[
 \beta^2_c:
 (R^2f_*\mathcal F)_c\otimes_{\mathcal O_{B,c}}\mathbb C(c)
 \xrightarrow{\sim} H^2(X_c,\mathcal F|_{X_c})
\]
is an isomorphism.  In particular, a nonzero target forces
$(R^2f_*\mathcal F)_c\ne0$.
\end{proposition}

\begin{proof}
B\u{a}nic\u{a}--St\u{a}n\u{a}\c{s}il\u{a}'s base-change theorem identifies
surjectivity of the degree-$q$ fibre map with right exactness of the
corresponding direct-image functor.  Their top-degree corollary states that
if $H^q(X_c,\mathcal F|_{X_c})=0$, then these equivalent conditions hold in
degree $q-1$, and the degree-$(q-1)$ base-change maps are isomorphisms
\cite[Chapter~III, \S3, Theorem~3.4 and Corollary~3.5,
pp.~116--121]{BanicaStanasila1976}.  Taking $q=3$ gives the displayed map.
If its target is nonzero, its source is nonzero; hence the stalk itself
cannot vanish.
\end{proof}

\begin{lemma}[automatic properness and flatness over the curve]
\label[lemma]{lem:automatic-proper-flat}
\Verified\ Let $X$ be a compact connected smooth complex manifold, let $B$
be a smooth complex curve, and let $f:X\to B$ be nonconstant and
holomorphic.  Then $f$ is proper and flat.  Every locally free
$\mathcal O_X$-module is $f$-flat, and every full scheme-theoretic fibre is
an effective Cartier divisor of pure codimension one.
\end{lemma}

\begin{proof}
The underlying continuous map is proper because $X$ is compact and $B$ is
Hausdorff.  Fix $x\in X$, put $b=f(x)$, and write
$A=\mathcal O_{B,b}$ and $R=\mathcal O_{X,x}$.  The ring $A$ is a discrete
valuation ring and $R$ is a regular local analytic ring, hence an integral
domain.  The homomorphism $A\to R$ is injective: if a nonzero germ on the
curve pulled back to zero near $x$, then a neighbourhood of $x$ would map
into its discrete zero set, so $f$ would be locally constant and therefore,
by connectedness and analytic continuation, globally constant.

Thus $R$ is torsion-free as an $A$-module.  A torsion-free module over a
discrete valuation ring is flat, so $f$ is flat at $x$; since $x$ was
arbitrary, $f$ is flat.  A locally free sheaf has stalk $R^r$, which is flat
over $A$ by the same argument.  Finally, if $t$ is a local parameter at
$b$, its pullback is a nonzero element of the domain $R$ and hence a
non-zero-divisor.  Therefore the full fibre is locally cut out by one
non-zero-divisor: it is an effective Cartier divisor and is pure of
codimension one.
\end{proof}

For a full fibre of complex dimension two and a coherent sheaf, $H^3=0$.
Proposition~\ref{prop:candidate-geometry} supplies the candidate geometry,
and Lemma~\ref{lem:automatic-proper-flat} makes the locally free sheaf
$\mathcal F_L=T_X\otimes L$ flat over the base.

\begin{corollary}[candidate base change and direct-image nonvanishing]
\label[corollary]{cor:candidate-r2}
\Verified\ For every $L\in\operatorname{Pic}(X)$, the canonical map
\[
 (R^2f_*\mathcal F_L)_{p_0}\otimes\mathbb C(p_0)
 \xrightarrow{\sim}H^2(W,i^*\mathcal F_L)
\]
is an isomorphism.  Consequently, for every
$L\in\operatorname{Pic}^0(X)$,
\[
 (R^2f_*\mathcal F_L)_{p_0}\ne0,
 \qquad\text{and hence}\qquad R^2f_*(T_X\otimes L)\ne0.
\]
If $r=0$, the same nonvanishing holds for every
$L\in\operatorname{Pic}(X)$.
\end{corollary}

\begin{proof}
The candidate is proper and $\mathcal F_L$ is base-flat, while the
two-dimensional fibre has $H^3(W,i^*\mathcal F_L)=0$.  Apply
Proposition~\ref{prop:top-degree-base-change}.  The nonzero target for
$L\in\operatorname{Pic}^0(X)$ is
Corollary~\ref{cor:candidate-fibrewise-h2}; the all-line-bundle extension is
the conditional final clause of that corollary.
\end{proof}

\subsection{Present evidentiary status}

\begin{itemize}
\item The reconstructed Sections~3--6 establish the compact smooth analytic
threefold, the proper map, the full reduced normal-crossings fibre, its
$\mathrm{dP}_6$ normalization, and the local triviality of $\omega_X$ near
$W$.  These claims have now been retraced independently from the frozen
source; they use none of the open integral-topology maps.
\item For every $L\in\operatorname{Pic}^0(X)$, the pullback of $A_L$ to the
normalization is trivial.  The conductor section, compact-Cartier duality,
and top-degree base change therefore prove
$R^2f_*(T_X\otimes L)\ne0$ for every such $L$.
\item Extending this conclusion to every $L\in\operatorname{Pic}(X)$ is a
separate question: for this method it is equivalent to $r=0$.  The claimed
$H^2(X;\mathbb Z)=0$ would suffice, but still depends on the open topology
bridges.
\item The conductor section invalidates the printed cotangent-sheaf
reduction as a formal inference on this non-normal fibre.  This does not by
itself decide whether CDP20, Theorem~2.2 has another proof or whether the
candidate satisfies all of that theorem's global hypotheses.
\end{itemize}

\subsection{Remaining verification and falsification programme}

\begin{enumerate}[label=\arabic*.]
\item \textbf{All-line-bundle restriction map.} Compute
$r=\eta^*i^*$ from explicit transition functions.  A line bundle with
$r(L)\ne0$ refutes the manuscript's universal conductor argument but does
not affect the proved $\operatorname{Pic}^0$ result.  A direct proof that
$r=0$ would extend the nonvanishing to all line bundles without borrowing
the unresolved topology.

\item \textbf{Global CDP hypotheses.} Establish independently the exact
topological vanishings, $a(X)=1$, and the holomorphic-algebraic-reduction
claim used to place this $X$ under CDP20, Theorem~2.2.  The analytic
nonvanishing above is not a substitute for those global premises.

\item \textbf{Literature comparison.} Reproduce the CDP p.~692 argument
with exactly the same fibre, twist, and definition of
$\widetilde\Omega^1_S$. Determine whether an omitted hypothesis restores
the inference or whether an independent proof bypasses it. Either outcome
must be recorded without converting a proof-gap diagnosis into a theorem
verdict.
\end{enumerate}

\begin{obligation}[CDP comparison]\label[obligation]{obl:cdp-comparison}
\Open\ Produce a line-by-line diagram comparing the candidate's
ambient, conormal, and intrinsic cotangent sequences with the exact sheaves
in the cited CDP reduction. Record separately: the validity of the printed
inference, the possibility of an independent repair, and the validity of
the candidate construction.
\end{obligation}

\begin{test}[decisive outcomes]
The proved $\operatorname{Pic}^0$ chain can be refuted only by locating an
error in one of its explicit inputs: the source construction of the full
fibre, the normalization--conductor descent, the $\mathrm{dP}_6$ Picard
injection, compact-Cartier duality, or top-degree base change.  The stronger
all-line-bundle chain fails exactly when the restriction map $r$ is nonzero.
Completing proof obligation~\ref{obl:cdp-comparison} would determine whether
the printed CDP reduction admits a repair, but would still leave the global
topology obligations
\hyperref[obl:C1]{C1}--\hyperref[obl:C3]{C3},
\hyperref[obl:V1]{V1}, \hyperref[obl:MV1]{MV1}, and
\hyperref[obl:L1]{L1}.
\end{test}
